\section*{Capítulo \ref{ch:g11:organic-skeletons} --- As moléculas orgânicas: esqueletos e nomes}

\begin{solution}{exo:g11:organic-skeletons:1}
As duas têm seis \omterm{def:g7:atoms-and-molecules:atom}{átomos} de carbono (extremidades e vértices) e
\ce{C6H14}: o hexano e o 2-metilpentano são \omterm{def:g11:organic-skeletons:isomer}{isômeros}.
\end{solution}

\begin{solution}{exo:g11:organic-skeletons:2}
Propano; hexano; 2-metilbutano.
\end{solution}

\begin{solution}{exo:g11:organic-skeletons:3}
\ce{CH3-CH2-CH3}; \chemfig{CH_3-CH(-[2]CH_3)-CH_3};
\chemfig{H-C(-[2]H)(-[6]H)-C(-[2]H)(-[6]H)-H}.
\end{solution}

\begin{solution}{exo:g11:organic-skeletons:4}
A primeira, o 3-metilpentano: 6 \omterm{def:g7:atoms-and-molecules:atom}{átomos} de carbono; as suas
extremidades \ce{CH3} (são três) levam 9 \omterm{def:g7:atoms-and-molecules:atom}{átomos} de hidrogênio, os seus
dois vértices \ce{CH2} levam 4 e o seu ponto de ramificação \ce{CH}
leva 1: \ce{C6H14}. O ciclo-hexano: 6 vértices, cada um \ce{CH2}:
\ce{C6H12}.
\end{solution}

\begin{solution}{exo:g11:organic-skeletons:5}
\ce{C8H18}. $2n + 2 = 20$ dá $n = 9$: o nonano, \ce{C9H20}.
\end{solution}

\begin{solution}{exo:g11:organic-skeletons:6}
\setchemfig{atom sep=1.6em}
3-metil-hexano \chemfig{-[:30]-[:-30]-[:30](-[:90])-[:-30]-[:30]-[:-30]},
\ce{C7H16}; 2,2-dimetilbutano \chemfig{-[:0](-[:90])(-[:-90])-[:0]-[:30]},
\ce{C6H14}; 3-etilpentano
\chemfig{-[:30]-[:-30](-[:-90]-[:-30])-[:30]-[:-30]}, \ce{C7H16}.
\end{solution}

\begin{solution}{exo:g11:organic-skeletons:7}
Pentano, 2-metilbutano e 2,2-dimetilpropano, desenhados na figura do
capítulo.
\end{solution}

\begin{solution}{exo:g11:organic-skeletons:8}
O ``2-etilbutano'' tem uma cadeia mais longa de cinco carbonos que passa
pelo grupo etil: é o 3-metilpentano. O ``4-metilpentano'' está numerado a
partir da extremidade errada: é o 2-metilpentano. O ``1-metilbutano'' é
simplesmente uma cadeia de cinco carbonos: o pentano.
\end{solution}

\begin{solution}{exo:g11:organic-skeletons:9}
2,2-dimetilpropano (\qty{9.5}{\celsius}) < 2-metilbutano
(\qty{27.8}{\celsius}) < pentano (\qty{36.1}{\celsius}). Os mesmos
\omterm{def:g7:atoms-and-molecules:atom}{átomos} estão organizados de modo cada vez mais compacto: menos contato
entre as \omterm{def:g7:atoms-and-molecules:molecule}{moléculas}, \omterm{def:g11:polarity-and-cohesion:van-der-waals}{interações de van der Waals} mais fracas, temperatura
de ebulição mais baixa.
\end{solution}

\begin{solution}{exo:g11:organic-skeletons:10}
Fechar a cadeia num anel cria mais uma ligação \ce{C-C}, que toma o
lugar de duas ligações \ce{C-H} (uma em cada extremidade). \omterm{def:g11:organic-skeletons:formulas}{Fórmulas
moleculares} diferentes: não são \omterm{def:g11:organic-skeletons:isomer}{isômeros}. O ciclo-hexano não é um \omterm{def:g11:organic-skeletons:alkane}{alcano}
(o seu esqueleto é um anel; a sua fórmula não é \ce{C_{n}H_{2n+2}}).
\end{solution}

\begin{solution}{exo:g11:organic-skeletons:11}
O butano e o 2-metilpropano (\ce{C4H10}). O propano é \ce{C3H8} e o
ciclobutano \ce{C4H8}.
\end{solution}

\begin{solution}{exo:g11:organic-skeletons:12}
\setchemfig{atom sep=1.5em}
\begin{center}
\small
\begin{tabular}{ccc}
\chemfig{-[:30]-[:-30]-[:30]-[:-30]-[:30]} &
\chemfig{-[:30](-[:90])-[:-30]-[:30]-[:-30]} &
\chemfig{-[:30]-[:-30](-[:-90])-[:30]-[:-30]} \\
hexano & 2-metilpentano & 3-metilpentano \\[6pt]
\chemfig{-[:0](-[:90])(-[:-90])-[:0]-[:30]} &
\chemfig{-[:30](-[:90])-[:-30](-[:-90])-[:30]} & \\
2,2-dimetilbutano & 2,3-dimetilbutano &
\end{tabular}
\end{center}
\end{solution}

\begin{solution}{exo:g11:organic-skeletons:13}
Dois zigue-zagues longos podem ficar lado a lado em todo o seu
comprimento, com muitos \omterm{def:g7:atoms-and-molecules:atom}{átomos} em contato; duas bolas se tocam só numa
pequena área. As \omterm{def:g11:polarity-and-cohesion:van-der-waals}{interações de van der Waals} atuam entre \omterm{def:g7:atoms-and-molecules:atom}{átomos} em
contato: elas são maiores entre as \omterm{def:g7:atoms-and-molecules:molecule}{moléculas} lineares, que precisam de
uma temperatura mais alta para se separar.
\end{solution}

\begin{solution}{exo:g11:organic-skeletons:14}
A cadeia mais longa tem seis carbonos; numerados a partir da esquerda,
os substituintes ficam em 2, 4, 4, e a partir da direita em 3, 3, 5; a
primeira diferença (2 contra 3) decide: 2,4,4-trimetil-hexano.
\end{solution}

\begin{solution}{exo:g11:organic-skeletons:15}
\setchemfig{atom sep=1.6em}
\chemfig{-[:0](-[:90])(-[:-90])-[:-30]-[:30](-[:90])-[:-30]}: uma cadeia
de cinco carbonos com dois grupos metil no carbono 2 e um no carbono 4,
oito \omterm{def:g7:atoms-and-molecules:atom}{átomos} de carbono ao todo; $2 \times 8 + 2 = 18$ \omterm{def:g7:atoms-and-molecules:atom}{átomos} de
hidrogênio, \ce{C8H18}, a fórmula do octano: eles são \omterm{def:g11:organic-skeletons:isomer}{isômeros}. \omterm{def:g11:organic-skeletons:formulas}{Fórmula
estrutural condensada}:
\chemfig{CH_3-C(-[2]CH_3)(-[6]CH_3)-CH_2-CH(-[2]CH_3)-CH_3}.
\end{solution}

\begin{solution}{pb:g11:organic-skeletons:1}
\setchemfig{atom sep=1.5em}
\textbf{1.} \ce{C4H10}: com $n = 4$, $2n + 2 = 10$.

\textbf{2.} \chemfig{H-C(-[2]H)(-[6]H)-C(-[2]H)(-[6]H)-C(-[2]H)(-[6]H)-C(-[2]H)(-[6]H)-H}.

\textbf{3.} \ce{CH3-CH2-CH2-CH3}; \chemfig{-[:30]-[:-30]-[:30]}.

\textbf{4.} \chemfig{H-C(-[2]H)(-[6]H)-C(-[6]H)(-[2,1.5]C(-[1]H)(-[3]H)(-[2]H))-C(-[2]H)(-[6]H)-H};
\chemfig{CH_3-CH(-[2]CH_3)-CH_3}; \chemfig{-[:30](-[:90])-[:-30]}. Ele
tem os mesmos \omterm{def:g7:atoms-and-molecules:atom}{átomos}, quatro carbonos e dez hidrogênios, ligados numa
ordem diferente.

\textbf{5.} Sim: \omterm{def:g11:organic-skeletons:isomer}{isômeros constitucionais} (esqueletos diferentes).

\textbf{6.} Os três: as suas temperaturas de ebulição estão abaixo de
\qty{20}{\celsius}.

\textbf{7.} A \qty{-5}{\celsius}, abaixo da sua temperatura de
ebulição, o butano continua líquido à pressão normal e libera pouco gás.

\textbf{8.} O propano (\qty{-42}{\celsius}) e o 2-metilpropano
(\qty{-11.7}{\celsius}) fervem abaixo das temperaturas de inverno: eles
ainda liberam bastante gás no frio.

\textbf{9.} O 2-metilpropano é ramificado, mais compacto: \omterm{def:g11:polarity-and-cohesion:van-der-waals}{interações de
van der Waals} mais fracas.

\textbf{10.} Com três \omterm{def:g7:atoms-and-molecules:atom}{átomos} de carbono, o único esqueleto possível é a
cadeia \ce{C-C-C}: qualquer carbono deslocado para outro lugar dá a
mesma cadeia.

\textbf{11.} 2-metilbutano.

\textbf{12.} A cadeia deve ser numerada a partir da extremidade mais
próxima do substituinte: 2, e não 3.

\textbf{13.} 2, 3 e 5.

\textbf{14.} Heptano, \chemfig{-[:30]-[:-30]-[:30]-[:-30]-[:30]-[:-30]}.

\textbf{15.} 2-metil-hexano \chemfig{-[:30](-[:90])-[:-30]-[:30]-[:-30]-[:30]}
e 3-metil-hexano \chemfig{-[:30]-[:-30](-[:-90])-[:30]-[:-30]-[:30]}.
O ``4-metil-hexano'' é o 3-metil-hexano numerado a partir da extremidade
errada.

\textbf{16.} 2,2-dimetilpentano, 2,3-dimetilpentano,
2,4-dimetilpentano e 3,3-dimetilpentano.

\textbf{17.} Sim, o 3-etilpentano. No carbono 2, o grupo etil formaria
uma cadeia de seis carbonos passando por ele: a \omterm{def:g7:atoms-and-molecules:molecule}{molécula} seria o
3-metil-hexano.

\textbf{18.} 2,2,3-trimetilbutano,
\chemfig{-[:0](-[:90])(-[:-90])-[:0](-[:90])-[:0]}.

\textbf{19.} $1 + 2 + 4 + 1 + 1 = 9$.

\textbf{20.} O heptano tem 9 \omterm{def:g11:organic-skeletons:isomer}{isômeros constitucionais}, contando ele
próprio.
\end{solution}
